\subsection{Basic Definitions and Hypotheses}\label{sec:conservation-basic-definitions}

We begin our study of conservation laws by computing their characteristics and giving conditions under which the systems are strictly hyperbolic.

\begin{Lemma}
\label{lem:conservation-characteristics}
A curve $t\mapsto\hat x(t)$ is a characteristic curve for the conservation law \eqref{eq:conservation-law-system} with solution $\bm{u}(x,t)$ if the matrix
\begin{equation}\label{eq:conservation-characteristic-matrix}
\hat x'(t)\bm{I}-\nabla\bm{f}(\bm{u}(\hat x(t),t))
\end{equation}
is singular. Furthermore, the system is strictly hyperbolic at a solution $\bm{u}$ if the eigenvalues of $\nabla\bm{f}(\bm{u})$ are real and distinct.
\end{Lemma}

The proof is left to the reader. All that is involved is interpreting the definition of a characteristic curve and strict hyperbolicity for a nonlinear system in the case where the curve is described by a graph rather than a level set. (Recall the comments about different representations for surfaces in Section~\ref{sec:cauchy-kovalevskaya}.)

Of course, the slopes of characteristic curves are nothing more than the eigenvalues of the matrix $\nabla\bm{f}(\bm{u})$. In light of this we introduce some notation describing eigenvalues and eigenvectors of $\nabla\bm{f}$. We assume that our system is strictly hyperbolic so that there are $n$ real distinct eigenvalues $\lambda_1(\bm{u})<\cdots<\lambda_n(\bm{u})$ with corresponding right and left eigenvectors $\bm{r}_k(\bm{u})$ and $\bm{l}_k(\bm{u})$ satisfying
\begin{equation}\label{eq:conservation-right-eigenvectors}
\nabla\bm{f}(\bm{u})\bm{r}_k(\bm{u})=\lambda_k(\bm{u})\bm{r}_k(\bm{u}),
\end{equation}

\begin{equation}\label{eq:conservation-left-eigenvectors}
\bm{l}_k(\bm{u})^\mathrm{T}\nabla\bm{f}(\bm{u})=\lambda_k(\bm{u})\bm{l}_k(\bm{u})^\mathrm{T}.
\end{equation}
Recall that since the eigenvectors are distinct, each of the sets of right and left eigenvectors $\{\bm{r}_1(\bm{u}),\allowbreak \ldots,\bm{r}_n(\bm{u})\}$ and $\{\bm{l}_1(\bm{u}),\allowbreak \ldots,\bm{l}_n(\bm{u})\}$ forms a basis for the state space.

We now define some functions on the state space, called \emph{Riemann invariants}, that are instrumental in finding solutions to problems with discontinuous initial conditions. These functions are defined locally in a neighborhood $U\subseteq\mathbb{R}^n$.

\begin{Definition}
\label{def:riemann-invariant}
A $k$-Riemann invariant is a smooth function $w:U\to\mathbb{R}$ such that for every $\bm{u}\in U$
\begin{equation}\label{eq:riemann-invariant-definition}
\bm{r}_k(\bm{u})\cdot\nabla w(\bm{u})=0.
\end{equation}

\end{Definition}

The following lemma gives an existence result for an appropriate system of Riemann invariants.

\begin{Lemma}
\label{lem:riemann-invariants-existence}
For every $\bar{\bm{u}}\in\mathbb{R}^n$ there is a neighborhood $\bar U\subset\mathbb{R}^n$ of $\bar{\bm{u}}$ on which there are $n-1$ $k$-Riemann invariants whose gradients are linearly independent at each point $\bm{u}\in\bar U$.
\end{Lemma}

\begin{Proof}
Let $S$ be a smooth surface through the point $\bar{\bm{u}}$, transversal to the vector $\bm{r}_k(\bar{\bm{u}})$. In a neighborhood of $\bar{\bm{u}}$, we now consider a system of the ODEs $\mathrm{d}\bm{u}/\mathrm{d}t=\bm{r}_k(\bm{u})$. Then $w(\bm{u})$ is a $k$-Riemann invariant if it is constant along every trajectory of this system of ODEs. Now every trajectory that passes through a sufficiently small neighborhood of $\bar{\bm{u}}$ intersects $S$ exactly once. The coordinates of this point of intersection (in a suitably chosen coordinate system on $S$) will serve as our Riemann invariants.
\end{Proof}

\begin{Example}{The $p$ system}
\label{ex:p-system-riemann-invariants}
We now consider the $p$ system
\begin{equation}\label{eq:p-system-vector-form}
\begin{pmatrix}w\\v\end{pmatrix}_t+\begin{pmatrix}-v\\p(w)\end{pmatrix}_x=\begin{pmatrix}0\\0\end{pmatrix}.
\end{equation}
Here we have
\begin{equation}\label{eq:p-system-flux-jacobian}
\nabla\bm{f}(w,v)=\begin{pmatrix}0&-1\\p'(w)&0\end{pmatrix}.
\end{equation}
To ensure strict hyperbolicity we assume
\begin{equation}\label{eq:p-system-strict-hyperbolicity}
p'<0.
\end{equation}
We now have eigenvalues $\lambda_1(w):=-\sqrt{-p'(w)}$ and $\lambda_2(w):=\sqrt{-p'(w)}$ with corresponding right eigenvectors
\begin{equation}\label{eq:p-system-first-eigenvector}
\bm{r}_1(w):=\begin{pmatrix}1\\\sqrt{-p'(w)}\end{pmatrix},
\end{equation}

\begin{equation}\label{eq:p-system-second-eigenvector}
\bm{r}_2(w):=\begin{pmatrix}1\\-\sqrt{-p'(w)}\end{pmatrix}.
\end{equation}
As indicated by the lemma above, there is one Riemann invariant corresponding to each eigenvalue; they are given as follows:
\begin{equation}\label{eq:p-system-first-riemann-invariant}
\rho_1(w,v):=v-\Psi(w),
\end{equation}

\begin{equation}\label{eq:p-system-second-riemann-invariant}
\rho_2(w,v):=v+\Psi(w),
\end{equation}
where
\begin{equation}\label{eq:p-system-riemann-primitive}
\Psi(w):=\int^w\sqrt{-p'(\xi)}\,\mathrm{d}\xi.
\end{equation}

\end{Example}

The relationship between the Riemann invariants and the characteristic curves for this system is given by the following result.

\begin{Theorem}
\label{thm:p-system-characteristic-invariants}
Let $(w(x,t),v(x,t))$ be a $C^1$ solution of the system given above. Then the Riemann invariant $\rho_i(w(x,t),v(x,t))$ is constant along characteristic curves satisfying $\hat x'(t)=-\lambda_i(w(\hat x(t),t))$.
\end{Theorem}

\begin{Proof}
We do only the calculation for $\rho_1$.
\[
\begin{aligned}
\dfrac{\mathrm{d}}{\mathrm{d}t}\rho_1(w(\hat x(t),t),v(\hat x(t),t))
&=v_x\hat x'+v_t-\sqrt{-p'(w)}(w_x\hat x'+w_t)\\
&=-\sqrt{-p'(w)}(w_t-v_x)+(v_t+p'(w)w_x)\\
&=0.
\end{aligned}
\]
The calculation for $\rho_2$ is identical.
\end{Proof}

One of the nice things about the $p$ system is that we can use the Riemann invariants as a convenient change of coordinates in state space; i.e., since the system is strictly hyperbolic, $\Psi'(w)=\sqrt{-p'(w)}>0$; hence $\Psi$ and the map $(w,v)\to(\rho_1,\rho_2)$ are invertible. If we rewrite the $p$ system in terms of $\rho_1,\rho_2$, we get the diagonal system
\begin{equation}\label{eq:p-system-diagonal-first}
\rho_{1,t}+\hat\lambda(\rho_1-\rho_2)\rho_{1,x}=0,
\end{equation}

\begin{equation}\label{eq:p-system-diagonal-second}
\rho_{2,t}-\hat\lambda(\rho_1-\rho_2)\rho_{2,x}=0.
\end{equation}
Here
\begin{equation}\label{eq:p-system-diagonal-speed}
\hat\lambda(s):=\sqrt{-p'\left(\Psi^{-1}\left(-\dfrac{s}{2}\right)\right)}.
\end{equation}


Both the Theorem~\ref{thm:p-system-characteristic-invariants} and the diagonalization procedure above can be generalized to any system of two strictly hyperbolic conservation laws. We can also use Riemann invariants to describe a hypothesis that often holds for systems of conservation laws coming from physics.

\begin{Definition}
\label{def:genuine-nonlinearity}
A system of conservation laws \eqref{eq:conservation-law-system} is said to be \textbf{genuinely nonlinear} in a region $D\subset\mathbb{R}^n$ if
\begin{equation}\label{eq:genuine-nonlinearity-definition}
\nabla\lambda_k\cdot\bm{r}_k\ne0,\quad k=1,2,\ldots,n
\end{equation}
in $D$.
\end{Definition}

\begin{Example}
\label{ex:scalar-genuine-nonlinearity}
In the case of a single conservation law \eqref{eq:scalar-conservation-law-quasilinear} we have $\lambda(u)=f'(u)$ and $r=1$, so $\nabla\lambda(u)\cdot r=f''(u)$. We refer to a function satisfying $f''>0$ ($<0$) as \emph{strongly convex} (\emph{concave}). In conservation laws, such a function is sometimes referred to as strictly convex (concave). This is (strictly speaking) incorrect. Thus, genuine nonlinearity is implied by either strong convexity or concavity of $f$. Strong convexity is often assumed for physical reasons. (Variational problems that represent the steady state of conservation laws are usually stated as minimization rather than maximization problems.)
\end{Example}

\begin{Example}
\label{ex:p-system-genuine-nonlinearity}
For the $p$ system we have
\begin{equation}\label{eq:p-system-genuine-nonlinearity}
\nabla\lambda_1\cdot\bm{r}_1=-\nabla\lambda_2\cdot\bm{r}_2=\dfrac{p''}{2\sqrt{-p'}}.
\end{equation}
Once again, strong convexity or strong concavity of $p$ is sufficient to ensure genuine nonlinearity. In typical applications in gas dynamics one assumes $p$ to be strongly convex.
\end{Example}

We should note that there are interesting physical problems that are not genuinely nonlinear. In particular, in the $p$ system the function $p$ is sometimes assumed to have an inflection point. We do not address such problems in detail in this book, but we should introduce the reader to the following terminology.

\begin{Definition}
\label{def:linear-degeneracy}
We say that the $k$th characteristic field is \textbf{linearly degenerate} at $\bm{u}$ if
\begin{equation}\label{eq:linear-degeneracy-definition}
\nabla\lambda_k(\bm{u})\cdot\bm{r}_k(\bm{u})=0.
\end{equation}

\end{Definition}
